\section{Introduction and Analysis Overview}
\label{sec:intro}
The Belle~II physics programme and detector are described in
Refs.~\cite{Kou:2018nap,Abe:2010gxa}.
\AnalysisTBD{physics motivation for $R(D)$, the current experimental and
theoretical status, and why an inclusive-tagging approach is complementary
to FEI-based measurements}

\subsection{Analysis signature}
\label{sec:intro:signature}
The signal-side $B$ is reconstructed as $B^{0} \to D^{-} \ell^{+}$ with
$D^{-} \to K^{+} \pi^{-} \pi^{-}$; the tag-side $B$ is represented by the
rest of the event rather than by an exclusive tag reconstruction
(Sec.~\ref{sec:reconstruction}).
Signal and normalisation yields are extracted from a two-dimensional binned
likelihood fit in $\mmiss$ and $\pdl$, performed simultaneously in a
$D$-mass signal region and a $D$-mass sideband (Sec.~\ref{sec:fit}).
The dominant experimental challenge is the generic-$B\bar{B}$ background,
which is treated in Sec.~\ref{sec:bbbar}.

\subsection{Analysis flow}
\label{sec:intro:flow}
\AnalysisTBD{one-page summary figure of the analysis flow}
